%!TEX root = thesis.tex
\chapter{Einleitung}
Die schnelle Entwicklung generativer KI-Systeme hat in den vergangenen Jahren zu erheblichen Fortschritten in der automatischen Erstellung von Texten sowie in der Verarbeitung komplexer Informationen geführt. Insbesondere sogenannte \textit{Large Language Models (LLMs)}, wie \textit{GPT-basierte} Modelle, \textit{Gemini} oder das auf Programmierung spezialisierte Modell \textit{GPT-5.3-Codex} zeigen bemerkenswerte Fähigkeiten in der Text- und Codeerstellung, der Problemlösung sowie in der Verarbeitung umfangreicher Wissensbestände \cite{brown2020languagemodelsfewshotlearners}.

Im Kontext der Softwareentwicklung eröffnen diese Modelle neue Möglichkeiten zur Unterstützung verschiedener Entwicklungsprozesse. LLMs können beispielsweise bei der Erstellung von Programmcodes der Analyse bestehender Systeme, der Erstellung von Dokumentationen oder der Entwicklung von Testfällen eingesetzt werden. Dadurch besteht die Möglichkeit, bestimmte Aufgaben im Softwareentwicklungsprozess effizienter zu gestalten und Entwickler bei komplexen Tätigkeiten zu unterstützen.
Gleichzeitig stehen viele Organisationen vor der Herausforderung, bestehende Softwaresysteme zu modernisieren. Insbesondere sogenannte Legacy-Systeme stellen häufig ein zentrales Problem dar, da sie über viele Jahre hinweg gewachsen sind und oft auf veralteten Technologien basieren. Zusätzlich fehlen in vielen Fällen aktuelle Dokumentationen oder eine klare Systemstruktur, was die Weiterentwicklung oder Migration solcher Systeme erheblich erschwert \cite{chen2021evaluatinglargelanguagemodels}.

Vor diesem Hintergrund gewinnt das sogenannte Softwarereengineering zunehmend an Bedeutung. Ziel des Reengineering, ist es, bestehende Systeme zu analysieren, zu verstehen und gegebenenfalls in eine modernere Architektur oder Implementierung zu überführen ohne dabei die ursprünglichen funktionalen Anforderungen zu verändern.
In diesem Zusammenhang stellt sich die Frage, inwieweit LLMs zur Unterstützung solcher Reengineering-Prozesse eingesetzt werden können. Aufgrund ihres Vermögens, Codes zu analysieren, natürliche Sprachbeschreibungen zu interpretieren und neue Implementierungen zu generieren, bieten LLMs die Möglichkeit, Teile der Systemmigration oder der Neuentwicklung bestehender Systeme zu automatisieren \cite{schlossarczyk2025ailegacy}.

Allerdings ist bislang nur begrenzt untersucht worden, in welchem Umfang LLM-basierte Ansätze tatsächlich in der Lage sind, bestehende Softwaresysteme unter Beibehaltung ihrer Anforderungen, Spezifikationen und Systemarchitekturen zuverlässig zu rekonstruieren oder neu zu implementieren. Insbesondere stellt sich die Frage, ob solche Modelle konsistente und funktional äquivalente Systeme erzeugen können.

Vor diesem Hintergrund untersucht diese Arbeit, inwiefern LLMs zur Unterstützung des Reengineerings bestehender Softwaresysteme eingesetzt werden können. Dabei wird analysiert, ob und in welchem Umfang LLMs dazu beitragen können, neue Systeme zu entwickeln, die die gleichen funktionalen Anforderungen und architektonischen Eigenschaften wie bestehende Systeme erfüllen.
\section{Problemstellung}
Viele Organisationen betreiben sogenannte Legacy-Systeme, die über viele Jahre hinweg entwickelt wurden und zentrale Geschäftsprozesse unterstützen \cite{sommerville2016software}. 
Die Modernisierung solcher Systeme stellt eine große Herausforderung 
im Software Engineering dar, da Legacy-Systeme häufig komplexe 
Abhängigkeiten und unzureichende Dokumentationen aufweisen 
\cite{seacord2003modernizing}. 
Ein möglicher Ansatz zur Modernisierung ist das sogenannte 
Softwarereengineering, bei dem bestehende Systeme analysiert 
und neu strukturiert werden, während ihre funktionalen 
Anforderungen erhalten bleiben \cite{chikofsky1990reverse}. 
Mit dem Aufkommen von LLMs eröffnen sich neue 
Möglichkeiten zur Unterstützung solcher Prozesse, beispielsweise 
durch automatisierte Codeanalyse oder Codegenerierung 
\cite{fan2023largelanguagemodelssoftware}.
Aktuelle Studien zeigen, dass Large Language Models in der Lage sind, Programm-codes auf Grundlage natürlicher Sprachbeschreibungen zu erzeugen \cite{chen2021evaluatinglargelanguagemodels}. Darüber hinaus werden LLMs zunehmend in verschiedenen Bereichen des Software Engineering eingesetzt,  beispielsweise zur Codeerstellung, Softwareanalyse oder Dokumentation \cite{LLMinSEreviewandvision}.

In dieser Arbeit wird untersucht: 
\textit{Inwieweit unterscheiden sich verschiedene Large Language 
Models hinsichtlich ihrer Fähigkeit, einen Programmcode zu erzeugen, 
der die funktionalen Anforderungen eines bestehenden Systems erfüllt?}


Das sogenannte Prompt-Engineering spielt eine wichtige Rolle bei der Steuerung der von LLMs generierten Ergebnisse, da die Formulierung der Eingabeaufforderung einen erheblichen Einfluss auf die Qualität der Ausgabe haben kann \cite{reynolds2021promptprogramminglargelanguage}.
Daher ist ein weiterer Aspekt dieser Arbeit die Untersuchung verschiedener Prompt-Strategien auf die Qualität und Funktionalität der Code-Erstellung. Als Anwendungsfall dient die Generierung eines webbasierten Buchungssystems.
